\chapter{On-line Analytical Processing Analysis}

The term OLAP refers to multidimensional analysis of large amounts of data. This is usually done with interactive tools with graphical interfaces that allow the user to make requests, which are automatically translated into SQL: for this purpose, SQL has been extended with some operators to support grouping and aggregation using analytic functions. This chapter will give an overview of OLAP systems, SQL data nalysis, and introduce reports.

\section{OLAP Systems Solutions}

An OLAP client provides graphical environments where the user can click or use drag-and-drop operations to provide input to the system. Experienced users can also write explicit SQL queries. The client can interact with the data manager using one of three possible solutions:
\begin{itemize}
    \item \textbf{DOLAP}: the client interacts with a local DOLAP (Desktop OLAP) system, which manages small amounts of data extracted from the main OLAP server, Data server, or operational DBMS. This is a good choice for users who need to work offline or move frequently (e.g., who travel a lot), or who do not need to perform complex queries. An example of DOLAP system are Excel Pivot Tables and Microsoft Power Pivot.
    \item \textbf{OLAP server}: the client interacts with an OLAP server, communicating via a dedicated API. The server hosts the multidimensional cube vision of a data warehouse; it can be of three types:
        \begin{itemize}
            \item \textbf{MOLAP} (Multidimensional OLAP), which stores both the data cube (taken from a relational Data server) and the aggregates in the extended cube in local memory, using a specialized multidimensional array-based data structure. MOLAP servers do not support SQL; instead, they support proprietary data cube query languages such as Microsoft's MDX. This solutions provides the best performances, but it is not suitable for large amounts of data.
            \item \textbf{ROLAP} (Relational OLAP), which stores both data and aggregates in the Data server. ROLAP servers may implement functionalities not supported by SQL.
            \item \textbf{HOLAP} (Hybrid OLAP), which combines MOLAP and ROLAP approaches: it splits data storage in a MOLAP server and a relational Data server. Splitting can be done in different ways. One method is to store the detailed data in the Data server, and keeping only precomputed, aggregate data in the MOLAP server; another method is to store the most recent aggregate data in the MOLAP server, to provide faster access, and older aggregate data in the Data server.
        \end{itemize}
        Data updates are usually done periodically in batches. Requests of the OLAP client to the server are formulated in SQL or other proprietary languages, and the result is typically returned in proprietary formats or in XML.
        \item \textbf{RDBMS} (Relational DBMS): the data warehouse is entirely stored in a relational Data server (i.e., a relational DBMS). All interactions with the OLAP client are directly expressed in SQL. An advantage of this solution is that is uses standard technology already available. In the past, it was considered to be too inadequate, both for performance limitations of RDBMSs, and lack of OLAP support in SQL. As modern RDBMS systems producers have added more and more specialized functionalities dedicated to OLAP applications (\textbf{OLAP-aware RDBMS}), this solution has become more and more popular.
\end{itemize}
All further discussions will assume the RDBMS solution.

\section{Reports}

This section is dedicated to reports, showing examples of how SQL can be used to generate the appropriate output for commonly asked business questions. All examples are based on the following table:
\begin{center}
    \textit{\textit{Sales}(Customer, Product, Brand, Date, City, Region, Area, Quantity, Revenue, Margin)}
\end{center}
Attributes don't have any null value.

\paragraph{Analytic SQL}
The SQL has been extended to allow the use of several analytic functions, also known as Windows Functions, to facilitate the development of complex data analysis. A query has the following structure:
\begin{lstlisting}[language=SQL, mathescape]
    SELECT Select_attr ($S_A$), Select_agg_fun ($S_{AF}$), Analytic_fun ($A_F$) OVER (
        [PARTITION BY <attr_list>]
        [ORDER BY <sort_attr_list>] [window_clause])
    FROM Fact_table ($F$) and Dimension_table (D)
    WHERE Where_cond ($W_C$)
    GROUP BY Grouping_attr ($G_A$)
    HAVING Having_cond ($H_C$) with Agg_fun ($H_{AF}$)
    ORDER BY Sort_attr ($O_A$);
\end{lstlisting}
A query written in analytic SQL is processed in this order: first, the operations in FROM, WHERE, GROUP BY, and HAVING clauses are performed; then, the analytic functions in the select are applied to the result of the previous step, producing a new result that adds new columns to the previous; finally, projection and sorting are applied to compute the final result.

Analytic functions are mostly used in the SELECT, with the \textbf{OVER} clause; they can either be applied to all the records, or separately to disjoint subsets obtained by partitioning the records by the value of an expression defined on the attributes of the record, using \textbf{PARTITION BY}. The partitioning operation is similar to that of a GROUP BY, but it does not output a single record for each group: instead, it produces as many records as it receives in input, extended with the new calculated attributes. If PARTITION BY is omitted, the set of records acts as a single group. Other analytic functions are specified in the GROUP BY clause, and define how to explicitly group records using some combinations of attributes, according to the chosen function.

\subsection{Simple Reports}
Simple reports are obtained from simple SQL queries. For example, take the following business question: "What was the total revenue of sales for the month of January 2025, by brand and by product?". The SQL query to answer this question and produce the correct report is:
\begin{lstlisting}[language=SQL]
    SELECT Brand, Product, SUM(Revenue) AS Total_Revenue
    FROM Sales
    WHERE YEAR(Date) = 2025 AND MONTH(Date) = 1
    GROUP BY Brand, Product
    ORDER BY Brand, Product;
\end{lstlisting}
Selection and projection are used to express slices and dices. Roll-ups and drill-downs are expressed using grouping: the former are obtained by removing attributes from the grouping clause, while the latter are obtained by adding attributes.

\paragraph{ROLLUP}
Suppose now we want to get a report with subtotals showing the total revenue from sales in the year 2025, by brand and by product, including the total revenue for each brand, and the overall total revenue. This can be done with standard SQL, although it requires three separate queries combined with \textit{UNION ALL}:
\begin{lstlisting}[language=SQL]
    SELECT Brand, Product, SUM(Revenue) AS Revenue
    FROM Sales
    WHERE YEAR(Date) = 2025
    GROUP BY Brand, Product

    UNION ALL
    
    SELECT Brand, NULL AS Product, SUM(Revenue) AS Revenue
    FROM Sales
    WHERE YEAR(Date) = 2025
    GROUP BY Brand

    UNION ALL

    SELECT NULL AS Brand, NULL AS Product, SUM(Revenue) AS Revenue
    FROM Sales
    WHERE YEAR(Date) = 2025
    ORDER BY Brand, Product;
\end{lstlisting}
The first query computes the revenue by brand and product, the second computes the subtotals for each brand, and the third computes the overall total. Fig. \ref{fig:report_rollup} shows an example of what the report would graphically look like.
\begin{figure}[H]
    \centering
    \begin{tabular}{lll}
        \hline
        & \textbf{Revenue by Brand and Product} & \\
        \hline
        \textbf{Brand} & \textbf{Product} & \textbf{Revenue} \\
        \hline
        B1 & P1 & 1000 \\
        & P2 & 1500 \\
        & P3 & 2000 \\
        \textbf{B1} & \textbf{Total} & \textbf{4500} \\
        B2 & P1 & 1200 \\
        & P3 & 1000 \\
        & P4 & 800 \\
        \textbf{B2} & \textbf{Total} & \textbf{3000} \\
        \textbf{Total} & & \textbf{7500} \\
        \hline
    \end{tabular}
    \caption{Report with subtotals for each brand.}
    \label{fig:report_rollup}
\end{figure}

Computing all these queries independently is inefficient: for this reason, many DBMSs support the \textbf{ROLLUP} clause, which automatically computes the totals and subtotals across the specified grouping attributes. Using ROLLUP, the previous query can be simplified as:
\begin{lstlisting}[language=SQL]
    SELECT Brand, Product, SUM(Revenue) AS Revenue
    FROM Sales
    WHERE YEAR(Date) = 2025
    GROUP BY ROLLUP(Brand, Product)
    ORDER BY Brand, Product;
\end{lstlisting}
The order in which attributes are specified in the ROLLUP clause is important, as it determines the order in which aggregates are computes. One way to interpret ROLLUP is as an operation that computes a path through the lattice of cuboids, starting from the most detailed one (i.e., the one with all attributes on the clause), and then removing one attribute at a time, in the reverse order in which they are specified in the ROLLUP clause, until no attribute is left. In this example, the order is: \textit{(Brand, Product)}, \textit{(Brand)}, \textit{()}.

\paragraph{CUBE}
Suppose we want to get a report similar to the one in Fig. \ref{fig:report_rollup}, but in addition, we want to add totals for each row and each column, meaning that we need to find the total revenue also for each product. Also in this case, a solution with standard SQL would require multiple queries combined, so it is better to use the \textbf{CUBE} clause included in many DBMSs. The query is in Fig. \ref{fig:report_cube}.
\begin{lstlisting}[language=SQL]
    SELECT Brand, Product, SUM(Revenue) AS Revenue
    FROM Sales
    WHERE YEAR(Date) = 2025
    GROUP BY CUBE(Brand, Product)
    ORDER BY Brand, Product;
\end{lstlisting}
Unlike roll-up, the order of the attributes does not matter. The report generated by this query is shown in Fig. \ref{fig:report_cube}.
\begin{figure}[H]
    \centering
    \begin{tabular}{lll}
        \hline
        & \textbf{Revenue by Brand and Product} & \\
        \hline
        \textbf{Brand} & \textbf{Product} & \textbf{Revenue} \\
        \hline
        B1 & P1 & 1000 \\
        & P2 & 1500 \\
        & P3 & 2000 \\
        \textbf{B1} & \textbf{Total} & \textbf{4500} \\
        B2 & P1 & 1200 \\
        & P3 & 1000 \\
        & P4 & 800 \\
        \textbf{B2} & \textbf{Total} & \textbf{3000} \\
        & \textbf{Total P1} & 2200 \\
        & \textbf{Total P2} & 1500 \\
        & \textbf{Total P3} & 3000 \\
        & \textbf{Total P4} & 800 \\
        \textbf{Total} & & \textbf{7500} \\
        \hline
    \end{tabular}
    \caption{Report with subtotals for each brand and each product.}
    \label{fig:report_cube}
\end{figure}

More than one ROLLUP or CUBE clause can be used in the same GROUP BY statement. For example, if we write:
\begin{lstlisting}[language=SQL]
    GROUP BY ROLLUP(A), ROLLUP(B, C)
\end{lstlisting}
the query will compute the cartesian product of both ROLLUPs; the first is $\{(\text{A}), ()\}$, and the second is $\{(\text{B,C}), (\text{B}), ()\}$, so their combination yields the following grouping sets:
\[
\{(\text{A, B, C}), (\text{A, B}), (\text{A}), (\text{B,C}), (\text{B}), ()\}
\]
This is also allowed for combinations of simple groupings with ROLLUPs and CUBEs, although in this case the result is constrained by the grouping attributes. For example, if we write:
\begin{lstlisting}[language=SQL]
    GROUP BY A, ROLLUP(B, C)
\end{lstlisting}
the result will be $\{(\text{A, B, C}), (\text{A, B}), (\text{A})\}$, because A must appear in all groupings. If we write:
\begin{lstlisting}[language=SQL]
    GROUP BY A, CUBE(B, C)
\end{lstlisting}
the result will be $\{(\text{A, B, C}), (\text{A, B}), (\text{A, C}), (\text{A})\}$. It is also possible to explicitly specify which attribute sets to group the data by with the \textbf{GROUPING SETS} clause:
\begin{lstlisting}[language=SQL]
    GROUP BY GROUPING SETS (
        (A, B, C),
        (B),
        ()
    )
\end{lstlisting}

\subsection{Moderately Difficult Reports}

Some commonly requested reports cannot be expressed easily, as they require comparisons. An example is: ``What are the total revenues in 2025 by brand and by product, with the percentage change from the previous year?''. One way to write the corresponding SQL query is to write two separate queries, one for the year 2025 and one for 2024, which individually find the total revenues grouping by brand and product. Then, if the same products are sold in both years, the final result can be found performing a natural join of the two results; otherwise, a full join is required. To write the full query, there are two options: either the subqueries are written in the FROM clause, or they are specified before the main one using the WITH clause. Using WITH, the query is:
\begin{lstlisting}[language=SQL, mathescape]
    WITH    Revenue25 AS(
            SELECT Brand, Product, SUM(Revenue) AS TotRevenue25 
            FROM Sales
            WHERE YEAR(Date) = 2025
            GROUP BY Brand, Product
            ),
            Revenue24 AS(
            SELECT Brand, Product, SUM(Revenue) AS TotRevenue25 
            FROM Sales
            WHERE YEAR(Date) = 2024
            GROUP BY Brand, Product
            )

    SELECT Revenue25.Brand AS Brand, Revenue25.Product AS Product,
    TotRevenue25, TotRevenue24,
    ROUND(100*(TotRevenue25 $-$ TotRevenue24)/TotRevenue25) AS Delta
    FROM Revenue24 FULL JOIN Revenue25 USING(Brand, Product)
    ORDER BY Brand, Product;
\end{lstlisting}
For simplicity, we omitted the handling of cases where the total revenue for a given product is NULL is a year it was not sold in; it would need an additional CASE statement.

Other business questions may require the use of OVER and PARTITION BY.

\subsection{Very Difficult Reports}

An example of a business question that needs a very difficult report to be answered is: ``What are the top 5 best products sold in a given region?''. This example motivates the introduction of an additonal operator in analytic SQL, called RANK.

\paragraph{RANK and DENSE\_RANK} The RANK function is used to sort the records in a set based on the value of an attribute or an expression (i.e., an aggregate function), assigning a position in a ranking to each record. The rank of a value is defined as 1 plus the number of other values that strictly precede it. If $k > 1$ values are equal, they are assigned the same rank, and the rank assigned to the next value is also increased by that $k$. For example, if the values are $(1, 2, 2, 3)$, their ranks are $(1, 2, 2, 4)$. The standard order is ascending, meaning that rank 1 is associated with the record with. the lowest value, but the ordering can be reversed. This function is specified in the SELECT clause as follows:
\begin{lstlisting}[language=SQL]
    <RankFun>()
    OVER ( [PARTITION BY <attr_list>]
        ORDER BY <sort_attr_list> ) [AS Ide]
\end{lstlisting}
Partitioning is optional. If it is omitted, all records are assumed to be part of the same partition set, so that each record is assigned a unique rank. If partitioning is used, ranking is done by grouping by the specified records and only within each group. The ORDER BY is mandatory, since ranking also requires the data to be sorted beforehand. DENSE\_RANK is a variant of RANK where the rank of a value is defined as 1 plus the number of distinct values that strictly precede it. Using the previous example, the values $(1, 2, 2, 3)$ would have ranks $(1, 2, 2, 3)$. If the collection of values contains no duplicates, RANK and DENSE\_RANK produce the same result. Other related functions include:
\begin{itemize}
    \item \textbf{PERCENT\_RANK}: defined as RANK - 1 divided by number of rows -1. It's a normalized version of RANK, since it always produces values in the range $[0, 1]$.
    \item \textbf{ROW\_NUMBER}: simply returns a sequential number for each record, starting from 1. 
    \item \textbf{CUME\_DIST}: it returns a value between 0 and 1 to each record according to the number of records lower than or equal to it. In other words, it calculates the empirical cumulative distribution of the values in the collection. Equal values are assigned to the same number.
    \item \textbf{NTILE($n$)}: returns the $n$-tile of the collection of records, i.e., a partition of $n$ groups each with the same number of records as the others (eventually plus or minus one if a perfect split is not possible), assigning to each record the number of the corresponding group.
\end{itemize}

\paragraph{LAG and LEAD} The LAG and LEAD functions are used to access data from rows that come before or after the current row in the result set, respectively. The syntax is:
\begin{lstlisting}[language=SQL]
    LAG(<attr>, <offset>, <default>)
    LEAD(<attr>, <offset>, <default>)
\end{lstlisting}
The first argument is the attribute from which the value is retrieved from. The second argument is the number of rows to skip above or below. The last argument is a default value to be used if the stored value in the specified row is empty (e.g., if the offset is larger than the number of rows available). Only the first argument is mandatory; the offset is otherwise set to 1, and the default is set to NULL.

\paragraph{Window Functions} Windows are used to make calculations that take into account a set of rows related to the current one. They can be used to compute cumulative, moving, and centered aggregates. The window size can be determined in a physical way with the \textbf{ROWS} option, or in a logical way using the \textbf{RANGE} option. Commonly, windows are defined over the date attribute. The window can include all records in the set on which it is defined, or only the current record. For example, to calculate a cumulative sum, the first record is fixed and the end of the window moves from first to last record. A moving average requires both window extremes to move. The syntax used to define a window is:
\begin{lstlisting}[language=SQL]
    <agg_fun>(<expr>)
    OVER (
        [PARTITION BY <attr_list>]
        [ORDER BY <sort_attr_list> [<ROWS or RANGE> <window_clause>]]
    ) [AS Ide]
\end{lstlisting}
Examples of window specifications are:
\begin{itemize}
    \item \textbf{ROWS UNBOUNDED PRECEDING}. The window begins with the first record of the partition and ends with the current record.
    \item \textbf{ROWS BETWEEN 5 PRECEDING AND 5 FOLLOWING}. The window includes all records that are within 5 rows before and after the current record.
    \item \textbf{RANGE BETWEEN INTERVAL 5 DAYS PRECEDING AND CURRENT ROW} (the date attribute is used). The window includes all records whose data falls within 5 days before the date of the current one.
\end{itemize}